\section{AHSME 1957}
        \begin{enumerate}
        	\item (\textit{AHSME 1957}) The number of distinct lines representing the altitudes, medians, and interior angle bisectors of a triangle that is isosceles, but not equilateral, is: 


 $\textbf{(A)}\ 9\qquad \textbf{(B)}\ 7\qquad \textbf{(C)}\ 6\qquad \textbf{(D)}\ 5\qquad \textbf{(E)}\ 3$

 \textbf{Answer: }B



\textbf{Solution}


                    \begin{figure}[H]
                        \centering
                        \includegraphics[width=0.5\textwidth]{Solution/AHSME_I_1957_1_Solution_Figure_1.png}
                    \end{figure}
                    

 As shown in the diagram above, all nine altitudes, medians, and interior angle bisectors are distinct, except for the three coinciding lines from the vertex opposite to the base. Thusly, there are $7$ distinct lines, so our answer is $\boxed{\textbf{(B)}}$, and we are done.

 


	\item (\textit{AHSME 1957}) In the equation $2x^2 - hx + 2k = 0$, the sum of the roots is $4$ and the product of the roots is $-3$. 
Then $h$ and $k$ have the values, respectively: 


 $\textbf{(A)}\ 8\text{ and }{-6} \qquad  \textbf{(B)}\ 4\text{ and }{-3}\qquad  \textbf{(C)}\ {-3}\text{ and }4\qquad \textbf{(D)}\ {-3}\text{ and }8\qquad \textbf{(E)}\ 8\text{ and }{-3}$

 \textbf{Answer: }E



\textbf{Solution}

Let the roots of the given equation be $r$ and $s$. Then, by \href{https://artofproblemsolving.com/wiki/index.php?title=Vieta\%27s\_Formulas}{Vieta's Formulas}, we have the following:\begin{align*}\frac{h}{2} = r+s = 4 &\rightarrow h = 8 \\\frac{2k}{2} = rs = -3 &\rightarrow k = -3 \\\end{align*}


 Thus, our answer is $\boxed{\textbf{(E) } 8 \text{ and } -3}$.

 


	\item (\textit{AHSME 1957}) The simplest form of $1 - \frac{1}{1 + \frac{a}{1 - a}}$ is: 


 $\textbf{(A)}\ {a}\text{ if }{a\not= 0} \qquad \textbf{(B)}\ 1\qquad  \textbf{(C)}\ {a}\text{ if }{a\not=-1}\qquad \textbf{(D)}\ {1-a}\text{ with not restriction on }{a}\qquad \textbf{(E)}\ {a}\text{ if }{a\not= 1}$

 \textbf{Answer: }E



\textbf{Solution}

We have $1 - \frac{1}{1 + \frac{a}{1 - a}} = 1 - \frac{1}{\frac{1}{1-a}} = 1 - \frac{1-a}{1} = a$ for almost all $a$. However, the first step is invalid when $a=1$, and each step is valid otherwise, so the answer is (E).

 


	\item (\textit{AHSME 1957}) The first step in finding the product $(3x + 2)(x - 5)$ by use of the distributive property in the form $a(b + c) = ab + ac$ is: 


 $\textbf{(A)}\ 3x^2 - 13x - 10 \qquad  \textbf{(B)}\ 3x(x - 5) + 2(x - 5)\qquad \\ \textbf{(C)}\ (3x+2)x+(3x+2)(-5)\qquad \textbf{(D)}\ 3x^2-17x-10\qquad \textbf{(E)}\ 3x^2+2x-15x-10$

 \textbf{Answer: }C



\textbf{Solution}

In order to find the product by using this form of the distributive property, we should express $(3x + 2)(x - 5)$ as $(3x + 2)(x) + (3x + 2)(-5)$. Thusly, $\boxed{\text{(C)}}$ is our answer, and we are done.

 


	\item (\textit{AHSME 1957}) Through the use of theorems on logarithms 

 \[\log{\frac{a}{b}} + \log{\frac{b}{c}} + \log{\frac{c}{d}} - \log{\frac{ay}{dx}}\]
can be reduced to: 


 $\textbf{(A)}\ \log{\frac{y}{x}}\qquad  \textbf{(B)}\ \log{\frac{x}{y}}\qquad  \textbf{(C)}\ 1\qquad \\ \textbf{(D)}\ 140x-24x^2+x^3\qquad \textbf{(E)}\ \text{none of these}$

 \textbf{Answer: }B



\textbf{Solution}

Using the properties $\log(x)+\log(y)=\log(xy)$ and $\log(x)-\log(y)=\log(x/y)$, we have
 \begin{align*} \log\frac ab+\log\frac bc+\log\frac cd-\log\frac{ay}{dx}&=\log\left(\frac ab\cdot\frac bc\cdot\frac cd\right)-\log \frac {ay}{dx} \\ &=\log \frac ad-\log\frac{ay}{dx} \\ &=\log\left(\frac{\frac ad}{\frac{ay}{dx}}\right) \\ &=\log \frac xy, \end{align*}
so the answer is $\boxed{\textbf{(B)} \log\frac xy}.$



	\item (\textit{AHSME 1957}) An open box is constructed by starting with a rectangular sheet of metal $10$ in. by $14$ in. 
and cutting a square of side $x$ inches from each corner. The resulting projections are folded up and the seams welded. 
The volume of the resulting box is: 


 $\textbf{(A)}\ 140x - 48x^2 + 4x^3 \qquad  \textbf{(B)}\ 140x + 48x^2 + 4x^3\qquad \\ \textbf{(C)}\ 140x+24x^2+x^3\qquad \textbf{(D)}\ 140x-24x^2+x^3\qquad \textbf{(E)}\ \text{none of these}$

 \textbf{Answer: }A



\textbf{Solution}

The resulting metal piece looks something like this where the white parts are squares of length $x$:

 
                    \begin{figure}[H]
                        \centering
                        \includegraphics[width=0.5\textwidth]{Solution/AHSME_I_1957_6_Solution_Figure_1.png}
                    \end{figure}
                    

 From here, try to visualize the rectangular prism coming together and realize the height is $x$, the length is $14-2x$, and the width is $10-2x$. Therefore, the volume is $x(14-2x)(10-2x)=x(4x^2-48x+40)= \boxed{\textbf{(A) } 140x - 48x^2 + 4x^3}$.

 


	\item (\textit{AHSME 1957}) The area of a circle inscribed in an equilateral triangle is $48\pi$. The perimeter of this triangle is: 


 $\textbf{(A)}\ 72\sqrt{3} \qquad \textbf{(B)}\ 48\sqrt{3}\qquad \textbf{(C)}\ 36\qquad \textbf{(D)}\ 24\qquad \textbf{(E)}\ 72$

 \textbf{Answer: }E



\textbf{Solution}


                    \begin{figure}[H]
                        \centering
                        \includegraphics[width=0.5\textwidth]{Solution/AHSME_I_1957_7_Solution_Figure_1.png}
                    \end{figure}
                    We can see that the radius of the circle is $4\sqrt{3}$. We know that the radius is $\frac{1}{3}$ of each median line of the triangle; each median line is therefore $12\sqrt{3}$. Since the median line completes a $30$-$60$-$90$ triangle, we can conclude that one of the sides of the triangle is $24$. Triple the side length and we get our answer, $\boxed{\textbf{(E)}72}$.

 


	\item (\textit{AHSME 1957}) The numbers $x,\,y,\,z$ are proportional to $2,\,3,\,5$. The sum of $x, y$, and $z$ is $100$. The number y is given by the equation $y = ax - 10$. Then a is: 


 $\textbf{(A)}\ 2 \qquad \textbf{(B)}\ \frac{3}{2}\qquad \textbf{(C)}\ 3\qquad \textbf{(D)}\ \frac{5}{2}\qquad \textbf{(E)}\ 4$

 \textbf{Answer: }A



\textbf{Solution}

In order to solve the problem, we first need to find each of the three variables. We can use the proportions the problem gives us to find the value of one part, and, by extension, the values of the variables (as $x$ would have $2$ parts, $y$ would have $3$, and $z$ would have $5$). One part, after some algebra, equals $10$, so $x$, $y$, and $z$ are $20$, $30$, and $50$, respectively.

 We can plug $x$ and $y$ into the equation given to us: $30 = 20a-10$, and then solve to get $a = \boxed{\textbf{(A)}2}$.

 


	\item (\textit{AHSME 1957}) The value of $x - y^{x - y}$ when $x = 2$ and $y = -2$ is: 


 $\textbf{(A)}\ -18 \qquad \textbf{(B)}\ -14\qquad \textbf{(C)}\ 14\qquad \textbf{(D)}\ 18\qquad \textbf{(E)}\ 256$

 \textbf{Answer: }B



\textbf{Solution}

Just plug in the numbers and follow the order of operations:
 \[2-(-2)^{2-(-2)}\]
 \[2-(-2)^4\]
 \[2-16\]
 \[\boxed{\textbf{(B) }-14}\]



	\item (\textit{AHSME 1957}) The graph of $y = 2x^2 + 4x + 3$ has its: 


 $\textbf{(A)}\ \text{lowest point at } {(-1,9)}\qquad   \textbf{(B)}\ \text{lowest point at } {(1,1)}\qquad \\  \textbf{(C)}\ \text{lowest point at }{(-1,1)}\qquad \textbf{(D)}\ \text{highest point at }{(-1,9)}\qquad\\  \textbf{(E)}\ \text{highest point at }{(-1,1)}$

 \textbf{Answer: }C



\textbf{Solution}

This graph generates a parabola, since the degree of $x$ is $2$. The $x-$ coordinate of the vertex of a parabola given by $ax^2 + bx + c$ is at $\frac{-b}{2a}$So, the vertex of this parabola is at 
 \[x = \frac{-4}{2(2)} = \frac{-4}{4} = -1\]Since the coefficient of $x^2$ is positive, at $x = -1$, the parabola is at its minimum. Substituting $x = -1$, we get
 \[y = 2(-1)^2 + 4(-1) + 3 \Rrightarrow y = 2 -4 + 3 \Rrightarrow y = 1\]So our answer is $\fbox{\textbf{(C)}}$.

 


 


	\item (\textit{AHSME 1957}) The angle formed by the hands of a clock at $2:15$ is: 


 $\textbf{(A)}\ 30^\circ \qquad \textbf{(B)}\ 27\frac{1}{2}^\circ\qquad  \textbf{(C)}\ 157\frac{1}{2}^\circ\qquad  \textbf{(D)}\ 172\frac{1}{2}^\circ\qquad \textbf{(E)}\ \text{none of these}$

 \textbf{Answer: }E



\textbf{Solution}

To find the angle of the clock, we first have to determine where the hands are. The time is $2.25$ hours, so the hour hand would be $67.5$ degrees clockwise. As the clock is a quarter of the way through the hour, the minute hand is $90$ degrees clockwise.

 Thus, we can say that the angle formed by the hands is $90 - 67.5 = 22\frac{1}{2}^\circ$.  This is answer  $\boxed{\textbf{(E)}}$.

 


	\item (\textit{AHSME 1957}) Comparing the numbers $10^{-49}$ and $2\cdot 10^{-50}$ we may say: 


 $\textbf{(A)}\ \text{the first exceeds the second by }{8\cdot 10^{-1}}\qquad\\  \textbf{(B)}\ \text{the first exceeds the second by }{2\cdot 10^{-1}}\qquad\\  \textbf{(C)}\ \text{the first exceeds the second by }{8\cdot 10^{-50}}\qquad\\  \textbf{(D)}\ \text{the second is five times the first}\qquad\\  \textbf{(E)}\ \text{the first exceeds the second by }{5}$

 \textbf{Answer: }C



\textbf{Solution}

By the rules of \href{https://artofproblemsolving.com/wiki/index.php?title=Exponentiation}{exponents}, $10^{-49}=10*10^{-50}$, which clearly exceeds $2*10^{-50}$. $10^{-49}-2*10^{-50}=10*10^{-50}-2*10^{-50}=8*10^{-50}$, so we choose answer choice $\boxed{\textbf{(C)}}$.

 


	\item (\textit{AHSME 1957}) A rational number between $\sqrt{2}$ and $\sqrt{3}$ is: 


 $\textbf{(A)}\ \frac{\sqrt{2} + \sqrt{3}}{2} \qquad  \textbf{(B)}\ \frac{\sqrt{2} \cdot \sqrt{3}}{2}\qquad  \textbf{(C)}\ 1.5\qquad \textbf{(D)}\ 1.8\qquad \textbf{(E)}\ 1.4$

 \textbf{Answer: }C



\textbf{Solution}

We see that $A$ and $B$ are both irrational, so we look at $C$, $D$, and $E$.Recall that $\sqrt 2$ is around $1.4$, and $\sqrt{3}$ is around $1.7$.The only number between these is $\boxed{\textbf{(C) }1.5}$.

 


	\item (\textit{AHSME 1957}) If $y = \sqrt{x^2 - 2x + 1} + \sqrt{x^2 + 2x + 1}$, then $y$ is: 


 $\textbf{(A)}\ 2x\qquad \textbf{(B)}\ 2(x+1)\qquad\textbf{(C)}\ 0\qquad\textbf{(D)}\ |x-1|+|x+1|\qquad\textbf{(E)}\ \text{none of these}$

 \textbf{Answer: }D



\textbf{Solution}

In order to solve the problem, we will use two properties, namely, that $(a+b)^2=a^2+2ab+b^2$ and $(a-b)^2=a^2-2ab+b^2$.

 We can use this to simplify the equation, as $y = \sqrt{x^2 - 2x + 1} + \sqrt{x^2 + 2x + 1}$ turns into $y = x + 1 + x - 1$. However, the square root function only allows for nonnegative inputs and only generates nonnegative outputs, so $y$ is $\boxed{\textbf{(D) }|x+1|+|x-1|}$.

 


	\item (\textit{AHSME 1957}) The table below shows the distance $s$ in feet a ball rolls down an inclined plane in $t$ seconds. 


 
 \[\begin{tabular}{|c|c|c|c|c|c|c|}\hline t & 0 & 1 & 2 & 3 & 4 & 5\\ \hline s & 0 & 10 & 40 & 90 & 160 & 250\\ \hline\end{tabular}\]


 The distance $s$ for $t = 2.5$ is: 


 $\textbf{(A)}\ 45\qquad \textbf{(B)}\ 62.5\qquad \textbf{(C)}\ 70\qquad \textbf{(D)}\ 75\qquad \textbf{(E)}\ 82.5$

 \textbf{Answer: }B



\textbf{Solution}

Looking at the pattern, we can determine that $t=10s^2$. Applying the relationship, we can see that $s = \boxed{\textbf{(B) }62.5}$ when $t=2.5$.

 


	\item (\textit{AHSME 1957}) Goldfish are sold at $15$ cents each. The rectangular coordinate graph showing the cost of $1$ to $12$ goldfish is: 


 $\textbf{(A)}\ \text{a straight line segment} \qquad \\  \textbf{(B)}\ \text{a set of horizontal parallel line segments}\qquad\\  \textbf{(C)}\ \text{a set of vertical parallel line segments}\qquad\\  \textbf{(D)}\ \text{a finite set of distinct points} \qquad\textbf{(E)}\ \text{a straight line}$

 \textbf{Answer: }D



\textbf{Solution}

$1$ goldfish costs $1 \cdot 15 = 15$ cents, $2$ goldfish cost $2 \cdot 15 = 30$ cents, and so on until $12$ fish. One may think that it is a straight line segment at first, but note that there cannot be a third of a goldfish sold (if so then that would be creepy). Thus our answer is $\boxed{\textbf{(B)}}$.

 


	\item (\textit{AHSME 1957}) A cube is made by soldering twelve $3$-inch lengths of wire properly at the vertices of the cube. 
If a fly alights at one of the vertices and then walks along the edges, the greatest distance it 
could travel before coming to any vertex a second time, without retracing any distance, is: 


 $\textbf{(A)}\ 24\text{ in.}\qquad \textbf{(B)}\ 12\text{ in.}\qquad \textbf{(C)}\ 30\text{ in.}\qquad  \textbf{(D)}\ 18\text{ in.}\qquad\textbf{(E)}\ 36\text{ in.}$

 \textbf{Answer: }A



\textbf{Solution}

To travel to an adjacent vertex, the fly must move $3$ inches. Furthermore, because a cube only has $8$ vertices, the fly can travel at most $8$ times before coming to a vertex for a second time (as a consequence of the \href{https://artofproblemsolving.com/wiki/index.php?title=Pigeonhole\_Principle}{Pigeonhole Principle}). Thus, the fly can travel at most for 24 inches. It is possible to reach this maximum by zig-zagging along the edges of the cube such that the fly walks over three of the four edges of each of the four side faces of the cube. Thus, our answer is $\boxed{\textbf{(A) }24 \text{ in.}}$.

 


	\item (\textit{AHSME 1957}) Circle $O$ has diameters $AB$ and $CD$ perpendicular to each other. $AM$ is any chord intersecting $CD$ at $P$. 
Then $AP\cdot AM$ is equal to: 


 
                    \begin{figure}[H]
                        \centering
                        \includegraphics[width=0.5\textwidth]{Problem/AHSME_I_1957_18_Problem_Figure_1.png}
                    \end{figure}
                    
$\textbf{(A)}\ AO\cdot OB \qquad \textbf{(B)}\ AO\cdot AB\qquad \\ \textbf{(C)}\ CP\cdot CD \qquad  \textbf{(D)}\ CP\cdot PD\qquad  \textbf{(E)}\ CO\cdot OP$

 \textbf{Answer: }B



\textbf{Solution 1}

Draw $MB$. Since $\angle AMB$ is inscribed on a diameter, $\angle AMB$ is $90^\circ$. By \href{https://artofproblemsolving.com/wiki/index.php?title=AA\_Similarity}{AA Similarity}, $\triangle APO \sim \triangle ABM$. Setting up ratios, we get $\frac{AP}{AO}=\frac{AB}{AM}$. Cross-multiplying, we get $AP\cdot AM = AO \cdot AB$, so the answer is $\boxed{\textbf{(B)}}$.

 



\textbf{Solution 2}

By \href{https://artofproblemsolving.com/wiki/index.php?title=Thales\%27\_Theorem}{Thales' Theorem}, $\measuredangle PMB = \measuredangle AMB = 90^{\circ}$. Because, from the problem, $\measuredangle POB = 90^{\circ}$ as well, $\measuredangle PMB + \measuredangle POB = 180^{\circ}$, so $PMBO$ is a \href{https://artofproblemsolving.com/wiki/index.php?title=Cyclic\_quadrilateral}{cyclic quadrilateral}. Thus, because $P$, $M$, $O$, and $B$ lie on a circle, we can use \href{https://artofproblemsolving.com/wiki/index.php?title=Power\_of\_a\_Point}{Power of a Point}. From this theorem, we get that $AP \cdot AM = AO \cdot AB$, which is answer choice $\boxed{\textbf{(B)}}$.

 


 


	\item (\textit{AHSME 1957}) The base of the decimal number system is ten, meaning, for example, that $123 = 1\cdot 10^2 + 2\cdot 10 + 3$. 
In the binary system, which has base two, the first five positive integers are $1,\,10,\,11,\,100,\,101$. 
The numeral $10011$ in the binary system would then be written in the decimal system as: 


 $\textbf{(A)}\ 19 \qquad \textbf{(B)}\ 40\qquad \textbf{(C)}\ 10011\qquad \textbf{(D)}\ 11\qquad \textbf{(E)}\ 7$

 \textbf{Answer: }A



\textbf{Solution}

Numbers in binary work similar to their decimal counterparts, where the multiplier associated with each place is multiplied by two every single place to the left. For example, $1111_2$ ($1111$ in base $2$) would equate to $1 * 2^3 + 1 * 2^2 + 1 * 2^1 + 1 * 2^0 = 8+4+2+1 = 15$.

 Using this same logic, $10011_2$ would be $1*2^4 + 1*2^1 + 1 * 2^0 = \boxed{\textbf{(A) }19}$.

 


	\item (\textit{AHSME 1957}) A man makes a trip by automobile at an average speed of 50 mph. He returns over the same route at an average speed of $45$ mph. 
His average speed for the entire trip is: 


 $\textbf{(A)}\ 47\frac{7}{19}\qquad  \textbf{(B)}\ 47\frac{1}{4}\qquad  \textbf{(C)}\ 47\frac{1}{2}\qquad  \textbf{(D)}\ 47\frac{11}{19}\qquad \textbf{(E)}\ \text{none of these}$

 \textbf{Answer: }A



\textbf{Solution}

Suppose the first half of the trip's distance is called $x$. Then the time for the first half is $\dfrac{x}{50},$ and the second half's time is $\dfrac{x}{45}$, so the total time is $\dfrac{x}{50}+\dfrac{x}{45}$, so the speed is $\dfrac{2x}{\frac{x}{50}+\frac{x}{45}}=\dfrac{2x}{\frac{19x}{450}}=2x \cdot \dfrac{450}{19x}=\dfrac{900}{19}=47\dfrac{7}{19}$, so the answer is $\boxed{\textbf{(A)}}$.

 


	\item (\textit{AHSME 1957}) Start with the theorem "If two angles of a triangle are equal, the triangle is isosceles," and the following four statements: 


 1. If two angles of a triangle are not equal, the triangle is not isosceles. 


 2. The base angles of an isosceles triangle are equal. 


 3. If a triangle is not isosceles, then two of its angles are not equal. 


 4. A necessary condition that two angles of a triangle be equal is that the triangle be isosceles. 


 Which combination of statements contains only those which are logically equivalent to the given theorem? 


 $\textbf{(A)}\ 1,\,2,\,3,\,4 \qquad \textbf{(B)}\ 1,\,2,\,3\qquad \textbf{(C)}\ 2,\,3,\,4\qquad \textbf{(D)}\ 1,\,2\qquad\textbf{(E)}\ 3,\,4$

 \textbf{Answer: }E



\textbf{Solution}

(1) is the inverse

 (2) is the converse

 (3) is the \href{https://artofproblemsolving.com/wiki/index.php?title=Contrapositive}{contrapositive}

 (4) is a restatement of the original theorem.

 Therefore, (3) and (4) are correct. $\boxed{\textbf{(E) } 3, 4}$



	\item (\textit{AHSME 1957}) If $\sqrt{x - 1} - \sqrt{x + 1} + 1 = 0$, then $4x$ equals: 


 $\textbf{(A)}\ 5 \qquad \textbf{(B)}\ 4\sqrt{-1}\qquad \textbf{(C)}\ 0\qquad \textbf{(D)}\ 1\frac{1}{4}\qquad \textbf{(E)}\ \text{no real value}$

 \textbf{Answer: }A



\textbf{Solution}

By repeatedly rearranging the equation and squaring both sides, we can solve for $x$:\begin{align*}\sqrt{x-1}-\sqrt{x+1}+1 &= 0 \\\sqrt{x-1}+1 &= \sqrt{x+1} \\x-1+2\sqrt{x-1}+1 &= x+1 \\2\sqrt{x-1} &= 1 \\\sqrt{x-1} &= \frac{1}{2} \\x-1 &= \frac{1}{4} \\x &= \frac{5}{4}\end{align*}
After checking for \href{https://artofproblemsolving.com/wiki/index.php?title=Extraneous\_solutions}{extraneous solutions}, we see that $x=\frac{5}{4}$ does indeed solve the equation. Thus, $4x=5$, and so our answer is $\boxed{\textbf{(A) }5}$.

 


	\item (\textit{AHSME 1957}) The graph of $x^2 + y = 10$ and the graph of $x + y = 10$ meet in two points. The distance between these two points is: 


 $\textbf{(A)}\ \text{less than 1} \qquad \textbf{(B)}\ 1\qquad \textbf{(C)}\ \sqrt{2}\qquad \textbf{(D)}\ 2\qquad\textbf{(E)}\ \text{more than 2}$

 \textbf{Answer: }C



\textbf{Solution}

We can merge the two equations to create $x^2+y=x+y$. Using either the quadratic equation or factoring, we get two solutions with $x$-coordinates $0$ and $1$.

 Plugging this into either of the original equations, we get $(0,10)$ and $(1,9)$. The \href{https://artofproblemsolving.com/wiki/index.php?title=Distance\_formula}{distance} between those two points is $\boxed{\textbf{(C) }\sqrt{2}}$.

 


	\item (\textit{AHSME 1957}) If the square of a number of two digits is decreased by the square of the number formed by reversing the digits, 
then the result is not always divisible by: 


 $\textbf{(A)}\ 9 \qquad  \textbf{(B)}\ \text{the product of the digits}\qquad  \textbf{(C)}\ \text{the sum of the digits}\qquad \textbf{(D)}\ \text{the difference of the digits}\qquad \textbf{(E)}\ 11$

 \textbf{Answer: }B



\textbf{Solution}

Let the first number have digits $\underline a$$\underline b$. Then, it equals $10a+b$, and the second number equals $10b+a$. Taking the difference of the squares of these two numbers and simplifying, we see that:

 \begin{align*}(10a+b)^2-(10b+a)^2 &= ((10a+b)-(10b+a))((10a+b)+(10b+a)) \\&= (9a-9b)(11a+11b) \\&= 9 \cdot 11(a-b)(a+b)\end{align*}
Thus, the result in question is always divisible by $9$, $11$, $(a-b)$, and $(a+b)$. This fact eliminates all of our options except $\fbox{\textbf{(B) }the product of the digits}$.

 


	\item (\textit{AHSME 1957}) The vertices of $\triangle PQR$ have coordinates as follows: $P(0,a),\,Q(b,0),\,R(c,d)$, where $a,\,b,\,c$ and $d$ are positive. 
The origin and point $R$ lie on opposite sides of $PQ$. The area of $\triangle PQR$ may be found from the expression: 


 $\textbf{(A)}\ \frac{ab + ac + bc + cd}{2} \qquad  \textbf{(B)}\ \frac{ac + bd - ab}{2}\qquad  \textbf{(C)}\ \frac{ab-ac-bd}{2}\qquad \textbf{(D)}\ \frac{ac+bd+ab}{2}\qquad \textbf{(E)}\ \frac{ac+bd-ab-cd}{2}$

 \textbf{Answer: }B



\textbf{Solution}


                    \begin{figure}[H]
                        \centering
                        \includegraphics[width=0.5\textwidth]{Solution/AHSME_I_1957_25_Solution_Figure_1.png}
                    \end{figure}
                    

 To solve this problem, we could use the \href{https://artofproblemsolving.com/wiki/index.php?title=Distance\_formula}{distance formula} to find the lengths of the sides and then \href{https://artofproblemsolving.com/wiki/index.php?title=Heron\%27s\_Formula}{Heron's Formula} to find the area of the triangle, but that solution seems messy and prone to mistakes with lots of square roots and polynomial expansions. Therefore, we look for a simpler, easier solution. Suppose $c=b$ and $a=d$. This makes $\triangle PQR$ half of a rectangle with side lengths $a$ and $b$, so it has an area of $\frac{ab}2$. Plugging in $c=b$ and $d=a$ into all of the answer choices, the only one which returns $\frac{ab}2$ is $\boxed{\textbf{(B) }\frac{ac + bd - ab}{2}}$.

 


	\item (\textit{AHSME 1957}) From a point within a triangle, line segments are drawn to the vertices. 
A necessary and sufficient condition that the three triangles thus formed have equal areas is that the point be: 


 $\textbf{(A)}\ \text{the center of the inscribed circle} \qquad \\  \textbf{(B)}\ \text{the center of the circumscribed circle}\qquad\\  \textbf{(C)}\ \text{such that the three angles formed at the point each be }{120^\circ}\qquad\\  \textbf{(D)}\ \text{the intersection of the altitudes of the triangle}\qquad\\  \textbf{(E)}\ \text{the intersection of the medians of the triangle}$

 \textbf{Answer: }E



\textbf{Solution}


                    \begin{figure}[H]
                        \centering
                        \includegraphics[width=0.5\textwidth]{Solution/AHSME_I_1957_26_Solution_Figure_1.png}
                    \end{figure}
                    

 Suppose the triangle is $\triangle ABC$ with the described point in its interior being $P$, as in the diagram. First, suppose that the smaller triangles have equal area (say $\frac{A}3$, where $A=[\triangle ABC]$). Then, by rearranging the area formula for a triangle, we see that the distance from $P$ to side $\overline{BC}$ is $\frac{2A}{3BC}$. Similarly, we can see that the distance from $A$ to $\overline{BC}$ is $\frac{2A}{BC}$. Thus, $P$ is $\frac1 3$ of the distance that $A$ is from $\overline{BC}$. We can use the same logic for points $B$ and $C$ and their respective opposite sides. The only point which is posiitioned in this way for all three points of the triangle is the \href{https://artofproblemsolving.com/wiki/index.php?title=Centroid}{centroid}, so it is necessary that $P$ be the intersection of the medians.

 Regarding the sufficiency of this condition, because the centroid is $\frac2 3$ of the way along each of the medians of the triangle, the three smaller triangles, with the same base and a third of the height of the large triangle, have one third of the area of the larger triangle. Thus, they all have equal areas, so it is sufficient that point $P$ is the centroid.

 Thus, our answer is $\fbox{\textbf{(E) }the intersection of the medians of the triangle}$.

 To disprove the other answers, try to draw counterexamples in extreme cases (so that it is obvious that a specific answer choice is incorrect).

 


	\item (\textit{AHSME 1957}) The sum of the reciprocals of the roots of the equation $x^2 + px + q = 0$ is: 


 $\textbf{(A)}\ -\frac{p}{q} \qquad \textbf{(B)}\ \frac{q}{p}\qquad \textbf{(C)}\ \frac{p}{q}\qquad \textbf{(D)}\ -\frac{q}{p}\qquad\textbf{(E)}\ pq$

 \textbf{Answer: }A



\textbf{Solution 1}

Let $f(x)=x^2+px+q$. Then, $x^2f(\frac1 x)$ equals $qx^2+px+1$ and has roots which are the reciprocals of those of $f(x)$. Thus, by \href{https://artofproblemsolving.com/wiki/index.php?title=Vieta\%27s\_Formulas}{Vieta's Formulas}, the sum of the roots of $x^2f(\frac1 x)$ (and thus the sum of the reciprocated roots of $f(x)$) is $\boxed{\textbf{(A) }\frac{-p}{q}}$.

 



\textbf{Solution 2}

One approach is to plug in some values for $p$ and $q$.

 We have $x^2-5x+6=0$.

 The roots are $x=2$ and $x=3$.

 The sum of the reciprocals of the roots is $\frac{1}{2}+\frac{1}{3}=\frac{5}{6}$.

 In this case, $p$ and $q$ are $-5$ and $6$.

 Thus, the answer is $\boxed{\textbf{(A) }\frac{-p}q}$.

 


	\item (\textit{AHSME 1957}) If $a$ and $b$ are positive and $a\not= 1,\,b\not= 1$, then the value of $b^{\log_b{a}}$ is: 


 $\textbf{(A)}\ \text{dependent upon }{b} \qquad \textbf{(B)}\ \text{dependent upon }{a}\qquad \textbf{(C)}\ \text{dependent upon }{a}\text{ and }{b}\qquad\textbf{(D)}\ \text{zero}\qquad\textbf{(E)}\ \text{one}$

 \textbf{Answer: }B



\textbf{Solution}

By the fact that logarithms are the \href{https://artofproblemsolving.com/wiki/index.php?title=Inverse\_of\_a\_function}{inverse} of exponentiation, it is clear that $b^{log_b{a}}=a$, so our answer is $\boxed{\textbf{(B)}\text{ dependent upon }a}$.

 


	\item (\textit{AHSME 1957}) The relation $x^2(x^2 - 1)\ge 0$ is true only for: 


 $\textbf{(A)}\ x \ge 1\qquad \textbf{(B)}\ - 1 \le x \le 1\qquad \textbf{(C)}\ x = 0,\, x = 1,\, x = - 1\qquad \\\textbf{(D)}\ x = 0,\, x\le-1,\, x\ge 1\qquad\textbf{(E)}\ x\ge 0$

 \textbf{Answer: }D



\textbf{Solution}

To solve this problem, think about the graph of $f(x)=x^2(x^2-1)=x^2(x-1)(x+1)$. The function equals zero only at the values $-1$, $0$, and $1$. Because the function is a quartic polynomial with a positive leading coefficient, it will go to positive infinity as $x$ tends to either positive or negative infinity. Thus, when $x$ is greatly negative, the function will be positive, and so the given inequality will hold. As $x$ gradually becomes more positive, it will eventually equal $-1$, when the function will equal zero. Thus, the $(x+1)$ term will switch to being positive, and so the whole function will become negative, where the inequality does not hold. $f(x)$ again reaches $0$ when $x=0$, but here the $x^2$ term does not change sign, so the function stays negative afterwards. Finally, when $x=1$, the function crosses the $x$-axis as the $(x-1)$ term changes sign, and the function goes off to positive infinity. Thus, the function is positive (and thus the given inequality will hold) when $x \leq -1$, $x=0$, and $x \geq 1$, which is answer choice $\fbox{\textbf{(D)}}$.

 


	\item (\textit{AHSME 1957}) The sum of the squares of the first n positive integers is given by the expression $\frac{n(n + c)(2n + k)}{6}$, 
if $c$ and $k$ are, respectively: 


 $\textbf{(A)}\ {1}\text{ and }{2} \qquad \textbf{(B)}\ {3}\text{ and }{5}\qquad \textbf{(C)}\ {2}\text{ and }{2}\qquad\textbf{(D)}\ {1}\text{ and }{1}\qquad\textbf{(E)}\ {2}\text{ and }{1}$

 \textbf{Answer: }D



\textbf{Solution}

When $n=1$, the value of the given expression must be $1^2=1$. Plugging in the values given by the answer choices, we see that the only option that returns $1$ when $n=1$ is $\boxed{\textbf{(D) }1 \text{ and } 1}$.

 


	\item (\textit{AHSME 1957}) A regular octagon is to be formed by cutting equal isosceles right triangles from the corners of a square. 
If the square has sides of one unit, the leg of each of the triangles has length: 


 $\textbf{(A)}\ \frac{2 + \sqrt{2}}{3} \qquad \textbf{(B)}\ \frac{2 - \sqrt{2}}{2}\qquad \textbf{(C)}\ \frac{1+\sqrt{2}}{2}\qquad\textbf{(D)}\ \frac{1+\sqrt{2}}{3}\qquad\textbf{(E)}\ \frac{2-\sqrt{2}}{3}$

 \textbf{Answer: }B



\textbf{Solution}


                    \begin{figure}[H]
                        \centering
                        \includegraphics[width=0.5\textwidth]{Solution/AHSME_I_1957_31_Solution_Figure_1.png}
                    \end{figure}
                    

 Let the side length of the regular octagon be $s$, and let the length of the legs of the isosceles right triangles be $x$. The triangles are $45-45-90$ triangles, so $s=x\sqrt2$. Because each side of the square is length $1$ and is composed of two legs of the triangles and one side of the octagon, $2x+s=1$. Substituting $s=x\sqrt2$ into this equation, we can now solve for $x$ to get our desired answer:\begin{align*}2x+s &= 1 \\2x+x\sqrt2 &= 1 \\x(2+\sqrt2) &= 1 \\x &= \frac{1}{2+\sqrt2} \cdot \frac{2-\sqrt2}{2-\sqrt2} \\x &= \frac{2-\sqrt2}{4-2} = \frac{2-\sqrt2}{2}\end{align*}
Thus, our answer is $\boxed{\textbf{(B) }\frac{2-\sqrt{2}}2}$.

 


	\item (\textit{AHSME 1957}) The largest of the following integers which divides each of the numbers of the sequence $1^5 - 1,\, 2^5 - 2,\, 3^5 - 3,\, \cdots, n^5 - n, \cdots$ is: 


 $\textbf{(A)}\ 1 \qquad \textbf{(B)}\ 60 \qquad \textbf{(C)}\ 15 \qquad \textbf{(D)}\ 120\qquad \textbf{(E)}\ 30$

 \textbf{Answer: }E



\textbf{Solution}

To begin, we look for a maximum value that this common divisor could have by looking at the smallest terms. $1^5-1=0$, but this gives us no information because all positive integers divide zero. However, $2^5-2=30$, so our answer cannot be greater than $30$ (because $30$ divided by an integer larger than $30$ is never an integer). This fact eliminates options (B) and (D).

 Now, we will factorize $n^5-n$ to discover some properties:\begin{align*}n^5-n &= n(n^4-1) \\&= n(n^2-1)(n^2+1) \\&= n(n-1)(n+1)(n^2+1)\end{align*}
One of $n$ and $n-1$ must be even, so $2$ divides $n^5-n$ for all $n$. Likewise, one of $n-1$, $n$, and $n+1$ must be divisible by $3$, so $3$ must divide $n^5-n$ as well. In the case of $5$, there are, in total, three cases:

 Case 1: $5$ divides either $n-1$, $n$, or $n+1$. Then, $5$ clearly must divide $n^5-n$.

 Case 2: $5$ divides $n+2$. In this case, $n+2 \equiv 0$ $\mod$ $5$, so $n \equiv 3 \mod 5$. Then, $n^2+1 \equiv 10 \equiv 0 \mod 5$, so $n^2-1$ (and, consequently, $n^5-n$) is divisible by $5$.

 Case 3: $5$ divides $n+3$. In this case, $n+3 \equiv 0 \mod 5$, so $n \equiv 2 \mod 5$. Thus, $n^2+1 \equiv 5 \equiv 0 \mod 5$, so $n^2-1$ (and, consequently, $n^5-n$) is divisible by $5$.

 Thus, 5 divides $n^5-n$ for all $n$.

 Because $2$, $3$, and $5$ all divide $n^5-n$ and are \href{https://artofproblemsolving.com/wiki/index.php?title=Relatively\_prime}{relatively prime}, $2 \cdot 3 \cdot 5 = 30$ also divides $n^5-n$. Thus, because $30$ is the largest of our remaining choices, we choose answer $\boxed{\textbf{(E) }30}$.

 


	\item (\textit{AHSME 1957}) If $9^{x + 2} = 240 + 9^x$, then the value of $x$ is: 


 $\textbf{(A)}\ 0.1 \qquad \textbf{(B)}\ 0.2\qquad \textbf{(C)}\ 0.3\qquad \textbf{(D)}\ 0.4\qquad \textbf{(E)}\ 0.5$

 \textbf{Answer: }E



\textbf{Solution}

$9^{x+2}$ can be rewritten as $9^x*9^2=9^x*81$, which means the equation can be rewritten as $81(9^x)=240+9^x$, or $80(9^x)=240$, or $9^x=3$. Therefore, $x=\boxed{\textbf{(E) }0.5}$.

 


	\item (\textit{AHSME 1957}) The points that satisfy the system $x + y = 1,\, x^2 + y^2 < 25$, constitute the following set: 


 $\textbf{(A)}\ \text{only two points} \qquad \\  \textbf{(B)}\ \text{an arc of a circle}\qquad \\  \textbf{(C)}\ \text{a straight line segment not including the end-points}\qquad\\  \textbf{(D)}\ \text{a straight line segment including the end-points}\qquad\\  \textbf{(E)}\ \text{a single point}$

 \textbf{Answer: }C



\textbf{Solution}


                    \begin{figure}[H]
                        \centering
                        \includegraphics[width=0.5\textwidth]{Solution/AHSME_I_1957_34_Solution_Figure_1.png}
                    \end{figure}
                    

 The line $x+y=1$ is in \href{https://artofproblemsolving.com/wiki/index.php?title=Standard\_form}{standard form}. Thus, it has a slope of $-\frac{1}{1}=-1$ and a $y$-intercept at $(0,1)$. It is graphed in blue in the diagram. All the points that satisfy this system must lie on this line. 

 The inequality $x^2+y^2<25$ is true for all points \href{https://artofproblemsolving.com/wiki/index.php?title=Strict\_inequality}{strictly} less than \href{https://artofproblemsolving.com/wiki/index.php?title=Distance\_formula}{distance} $\sqrt{25}=5$ from the origin. This area is a \href{https://artofproblemsolving.com/wiki/index.php?title=Disk}{disk} with radius $5$, whose open boundary is represented in red in the diagram.

 Because the $y$-intercept of the line, $(0,1)$, is inside the disk, it cannot be tangent to the disk. Thus, there must be infiniely many points along the line which are within the disk and thereby satisfy the system. However, because the inequality is strict, the segment of solutions does not include the endpoints, so our answer is $\fbox{\textbf{(C)} a straight line segment not including the end-points}$.

 


	\item (\textit{AHSME 1957}) Side $AC$ of right triangle $ABC$ is divided into $8$ equal parts. Seven line segments parallel to $BC$ are drawn to $AB$ from the points of division. 
If $BC = 10$, then the sum of the lengths of the seven line segments: 


 $\textbf{(A)}\ \text{cannot be found from the given information} \qquad \textbf{(B)}\ \text{is }{33}\qquad \textbf{(C)}\ \text{is }{34}\qquad\textbf{(D)}\ \text{is }{35}\qquad\textbf{(E)}\ \text{is }{45}$

 \textbf{Answer: }D



\textbf{Solution}


                    \begin{figure}[H]
                        \centering
                        \includegraphics[width=0.5\textwidth]{Solution/AHSME_I_1957_35_Solution_Figure_1.png}
                    \end{figure}
                    

 Because the location of the right angle within the triangle is not specified, answer choice (A) may be tempting. However, as we will come to see later, this ambiguity is irrelevant to the problem. Because we have a series of parallel lines, we have many smaller triangles which are \href{https://artofproblemsolving.com/wiki/index.php?title=Similarity\_(geometry)}{similar} to the larger triangle. Because the smallest triangle at the top has its base $\frac1 8$ of the way along $\overline{AC}$, its base must be $\frac1 8$ of that of the larger triangle, so it must have length $\frac1 8 \cdot 10$. We can repeat this for each of the seven segments and thus find our desired sum (by using the formula for \href{https://artofproblemsolving.com/wiki/index.php?title=Triangular\_numbers}{triangular numbers}): \begin{align*}&10 \cdot \frac1 8 + 10 \cdot \frac2 8 + ... + 10 \cdot \frac7 8 \\= &\frac{10}8(1+2+...+7) \\= &\frac{10}8(\frac{7 \cdot 8}2) \\= &10 \cdot \frac7 2 \\= &35\end{align*}
Thus, our answer is $\boxed{\textbf{(D)} \text{ is } 35}$.Note that this answer holds for any $\triangle ABC$ such that $BC=10$, so the aforementioned ambiguity is irrelevant.

 


	\item (\textit{AHSME 1957}) If $x + y = 1$, then the largest value of $xy$ is: 


 $\textbf{(A)}\ 1\qquad \textbf{(B)}\ 0.5\qquad \textbf{(C)}\ \text{an irrational number about }{0.4}\qquad \textbf{(D)}\ 0.25\qquad\textbf{(E)}\ 0$

 \textbf{Answer: }D



\textbf{Solution}

By \href{https://artofproblemsolving.com/wiki/index.php?title=AM-GM}{AM-GM}, we have
 \[\frac{x+y}{2} \geq \sqrt{xy}\]Substituting, we have
 \[\frac{1}{2} \geq \sqrt {xy}\]
 \[\frac{1}{4} \geq xy\]Equality occurs when $x = y = \boxed{\textbf{(D) }\frac12}$.

 


	\item (\textit{AHSME 1957}) In right triangle $ABC, BC = 5, AC = 12$, and $AM = x; \overline{MN} \perp \overline{AC}, \overline{NP} \perp \overline{BC}$; 
$N$ is on $AB$. If $y = MN + NP$, one-half the perimeter of rectangle $MCPN$, then: 


 
                    \begin{figure}[H]
                        \centering
                        \includegraphics[width=0.5\textwidth]{Problem/AHSME_I_1957_37_Problem_Figure_1.png}
                    \end{figure}
                    
$\textbf{(A)}\ y = \frac {1}{2}(5 + 12) \qquad  \textbf{(B)}\ y = \frac {5x}{12} + \frac {12}{5}\qquad  \textbf{(C)}\ y =\frac{144-7x}{12}\qquad \\ \textbf{(D)}\ y = 12\qquad \qquad\quad\,\,  \textbf{(E)}\ y = \frac {5x}{12} + 6$

 \textbf{Answer: }C



\textbf{Solution}

Because $AC=12$ and $AM=x$, $MC=12-x$. Let $MN=z$. Then, because $MNPC$ is a rectangle, $NP=12-x$ and $PC=z$, and so $BP=5-z$. By \href{https://artofproblemsolving.com/wiki/index.php?title=AA\_similarity}{AA similarity}, $\triangle AMN \sim \triangle NPB$. From this similarity, we can solve the following proportion for $z$:\begin{align*}\frac{BP}{PN} &= \frac{NM}{AM} \\\frac{5-z}{12-x} &= \frac z x \\5x-xz &= 12z-xz \\5x &= 12z \\z &= \frac{5x}{12}\end{align*}
Because $y=MN+NP=z+12-x$, we can now substitute for $z$ and find $y$ in terms of $x$:\begin{align*}y &= z+12-x \\&= \frac{5x}{12}-x+12 \\&= \frac{-7x}{12}+\frac{144}{12} \\&= \frac{144-7x}{12}\end{align*}
Thus, our answer is $\boxed{\textbf{(C) }\frac{144-7x}{12}}$.

 


	\item (\textit{AHSME 1957}) From a two-digit number $N$ we subtract the number with the digits reversed and find that the result is a positive perfect cube. Then: 


 $\textbf{(A)}\ {N}\text{ cannot end in 5}\qquad\\  \textbf{(B)}\ {N}\text{ can end in any digit other than 5}\qquad \\  \textbf{(C)}\ {N}\text{ does not exist}\qquad\\  \textbf{(D)}\ \text{there are exactly 7 values for }{N}\qquad\\  \textbf{(E)}\ \text{there are exactly 10 values for }{N}$

 \textbf{Answer: }D



\textbf{Solution}

The number $N$ can be written as $10a+b$ with $a$ and $b$ representing the digits. The number $N$ with its digits reversed is $10b+a$. Since the problem asks for a positive number as the difference of these two numbers, than $a>b$. Writing this out, we get $10a+b-(10b+a)=9a-9b=9(a-b)$. Therefore, the difference must be a multiple of $9$, and the only perfect cube with less than $3$ digits and is multiple of $9$ is $3^3=27$. Also, that means $a-b=3$, and there are $7$ possibilities of that, so our answer is $\boxed{\textbf{(D) } \text{There are exactly } 7 \text{ values for } N}$.

 


	\item (\textit{AHSME 1957}) Two men set out at the same time to walk towards each other from $M$ and $N$, $72$ miles apart. 
The first man walks at the rate of $4$ mph. The second man walks $2$ miles the first hour, $2\frac {1}{2}$ 
miles the second hour, $3$ miles the third hour, and so on in arithmetic progression. Then the men will meet: 


 $\textbf{(A)}\ \text{in 7 hours} \qquad  \textbf{(B)}\ \text{in }{8\frac {1}{4}}\text{ hours}\qquad  \textbf{(C)}\ \text{nearer }{M}\text{ than }{N}\qquad\\  \textbf{(D)}\ \text{nearer }{N}\text{ than }{M}\qquad \textbf{(E)}\ \text{midway between }{M}\text{ and }{N}$

 \textbf{Answer: }E



\textbf{Solution}

Let $t$ be the time it takes (in hours) for the men to meet. Then, the distance the first man will have travelled is $4t$. If $t$ is an integer, the second man will have travelled $\frac{2+2+0.5(t-1)}{2} \cdot t$, the sum of the first $t$ terms of the \href{https://artofproblemsolving.com/wiki/index.php?title=Arithmetic\_progression}{arithmetic progression}. If $t$ is not an integer, then we should have a good enough approximation to choose an answer choice (given the relatively large distance of $72$ miles compared to their speeds). Because the combined distance that the men travel must be $72$ miles, we can now solve for $t$ in the following formula:\begin{align*}4t+\frac{2+2+0.5(t-1)}{2} \cdot t &= 72 \\4t+\frac{4+0.5t-0.5}{2} \cdot t &= 72 \\8t+3.5t+0.5t^2 &= 144 \\t^2+23t &= 288 \\t^2+23t-288 &= 0 \\(t+32)(t-9) &= 0\end{align*}
Because $t>0$, we are left with $t=9$. At $t=9$, the first man will have travelled $36$ miles, which is half of te distance between $M$ and $N$. Thus, our answer is $\boxed{\textbf{(E)} \text{ midway between } M \text{ and } N}$.

 


	\item (\textit{AHSME 1957}) If the parabola $y = -x^2 + bx - 8$ has its vertex on the $x$-axis, then $b$ must be: 


 $\textbf{(A)}\ \text{a positive integer}\qquad \\  \textbf{(B)}\ \text{a positive or a negative rational number}\qquad \\  \textbf{(C)}\ \text{a positive rational number}\qquad\\  \textbf{(D)}\ \text{a positive or a negative irrational number}\qquad\\  \textbf{(E)}\ \text{a negative irrational number}$

 \textbf{Answer: }D



\textbf{Solution 1}

Note that if $y=-x^2+bx-8$ has its vertex on the $x$-axis, then $-y=x^2-bx+8$ will have its vertex on the x-axis as well. To find the location of the vertex of the parabola, we desire to put it in vertex form, where $-y=(x-h)^2+k$, and $(h,k)$ is the location of the vertex. However, we know that $k=0$, because the vertex is on the $x$-axis. Thus, we know that $x^2-bx+8$ must be the \href{https://artofproblemsolving.com/wiki/index.php?title=Perfect\_square\#Perfect\_Square\_Trinomials}{square of a linear term}. Thus, $b=\pm 2 \cdot 1 \cdot \sqrt8 =\pm 4\sqrt2$, which are both irrational. Thus, our answer is $\boxed{\textbf{(D)}\text{ a positive or a negative irrational number}}$.

 



\textbf{Solution 2}

We know that if a parabola is given by $ax^2+bx+c$, then the $x$-value of the vertex is $\frac{-b}{2a}$ (this fact can be proven with the \href{https://artofproblemsolving.com/wiki/index.php?title=Quadratic\_formula}{quadratic formula} and also \href{https://artofproblemsolving.com/wiki/index.php?title=Derivative}{derivatives}). Because, in this case, $a=-1$, $x=\frac{-b}{2a}=\frac b 2$. Thus, at $x=\frac b 2$, the parabola should have a $y$-value of $0$. Therefore, we have the following equation that we can solve for $b$:\begin{align*}y|_{x=\frac b 2} = -(\frac b 2)^2-b \cdot (\frac b 2)-8 &= 0\\-\frac{b^2}4+\frac{b^2}2 &= 8 \\\frac{b^2}4 &= 8 \\b^2 &= 32 \\b &= \pm \sqrt{32} = \pm 4\sqrt2\end{align*}
Because both of these results are irrational with one positive and one negtaive solution, we choose answer $\fbox{\textbf{(D)} a positive or a negative irrational number}$.

 


	\item (\textit{AHSME 1957}) Given the system of equations 

 \[ax + (a - 1)y = 1\]

 \[(a + 1)x - ay = 1\]
For which one of the following values of $a$ is there no solution $x$ and $y$? 


 $\textbf{(A)}\ 1\qquad  \textbf{(B)}\ 0\qquad  \textbf{(C)}\ - 1\qquad  \textbf{(D)}\ \frac {\pm \sqrt {2}}{2}\qquad  \textbf{(E)}\ \pm\sqrt{2}$

 \textbf{Answer: }D



\textbf{Solution}

Notice that the two equations are lines in \href{https://artofproblemsolving.com/wiki/index.php?title=Standard\_form}{standard form}, so the system will have no solutions if these lines never intersect. For them never to intersect, they must have the same slope and not be the exact same line. We are sure that they are not the same line, because they have different coefficients and have the same number, $1$, on their right hand sides. In the case that one of the lines has an undefined slope (where it is of the form $x=k$, where $k$ is some constant), then $a$ is either $1$ or $0$. In either case, the other line does not have an undefined slope, so the lines eventually intersect and there must be a solution to the system of equations. Thus, we need not worry about this case. Because the lines are in standard form, we can equate the two expressions for their slope as follows:\begin{align*}-\frac{a-1}{a} &= - \frac{-a}{a+1} \\-(a^2-1) &= a^2 \\2a^2 &= 1 \\a^2 &= \frac 1 2 \\a &= \pm \frac 1{\sqrt2} = \pm \frac{\sqrt2}2\end{align*}
Thus, our answer is $\boxed{\textbf{(D) }\frac{\pm \sqrt2}2}$.

 


	\item (\textit{AHSME 1957}) If $S = i^n + i^{-n}$, where $i = \sqrt{-1}$ and $n$ is an integer, then the total number of possible distinct values for $S$ is: 


 $\textbf{(A)}\ 1\qquad \textbf{(B)}\ 2\qquad \textbf{(C)}\ 3\qquad \textbf{(D)}\ 4\qquad \textbf{(E)}\ \text{more than 4}$

 \textbf{Answer: }C



\textbf{Solution 1}

We first use the fact that $i^{-n}=\frac1{i^n}=\left(\frac1i\right)^n=(-i)^n$. Note that $i^4=1$ and $(-i)^4=1$, so $i^n$ and $(-i)^n$ are periodic with periods at most 4. Therefore, it suffices to check for $n=0,1,2,3$.

 
For $n=0$, we have $i^0+(-i)^0=1+1=2$.

 For $n=1$, we have $i^1+(-i)^1=i-i=0$.

 For $n=2$, we have $i^2+(-i)^2=-1-1=-2$.

 For $n=3$, we have $i^3+(-i)^3=-i+i=0$.

 Hence, the answer is $\boxed{\textbf{(C)}\ 3}$.

 



\textbf{Solution 2}

Notice that the powers of $i$ cycle in cycles of 4. So let's see if $S$ is periodic.

 For $n=0$: we have $2$. 

 For $n=1$: we have $0$.

 For $n=2$: we have $-2$.

 For $n=3$: we have $0$.

 For $n=4$: we have $2$ again. Well, it can be seen that $S$ cycles in periods of 4. Select $\fbox{\textbf{(C)}}$.

 


	\item (\textit{AHSME 1957}) We define a lattice point as a point whose coordinates are integers, zero admitted. 
Then the number of lattice points on the boundary and inside the region bounded by the $x$-axis, 
the line $x = 4$, and the parabola $y = x^2$ is: 


 $\textbf{(A)}\ 24 \qquad  \textbf{(B)}\ 35\qquad  \textbf{(C)}\ 34\qquad  \textbf{(D)}\ 30\qquad  \textbf{(E)}\ \infty$

 \textbf{Answer: }B



\textbf{Solution 1}


                    \begin{figure}[H]
                        \centering
                        \includegraphics[width=0.5\textwidth]{Solution/AHSME_I_1957_43_Solution_Figure_1.png}
                    \end{figure}
                    

 We want to find the number of lattice points in and on the boundary of the shaded region in the diagram. To do this, we will look at the integer values of $x$ from $0$ to $4$. At a given value of $x$, the amount of lattice points in the region is $x^2+1$, because all of the integers from $0$ up to and including $x^2$ are in the region. Thus, evaluating this expression at at $x=0,1,2,3,$ and $4$ and adding the results together, we see that the number of lattice points is $1+2+5+10+17=35$, so our answer is $\boxed{\textbf{(B) }35}$.

 



\textbf{Solution 2 (not as fast, but cool theorems)}

Let us construct a simple polygon with vertices at $(0,0), (1,1), (2,4), (3,9), (4,16),$ and $(4,0)$. Notice that this polygon overestimates the area of the shaded region in the diagram (because the parabola has \href{https://artofproblemsolving.com/wiki/index.php?title=Convex\_function}{positive concavity}) yet does not have any more lattice points, because it matches the parabola at integral values of $x$, and lattice points can only exist when $x$ is an integer. After using the \href{https://artofproblemsolving.com/wiki/index.php?title=Shoelace\_Theorem}{Shoelace Theorem}, we see that the area of the polygon is $\frac{114-70}2=22$. Let $B$ be the number of lattice points on the boundary of the polygon, and let $I$ be the number of lattice points inside the polygon. We know that $B=5+17+5-3=24$ by counting the lattice points on the parabola, on the right edge, and on the bottom edge while taking care to subtract out the three points we counted twice (namely, $(0,0), (4,16),$ and $(4,0)$). Now, recalling that the polygon has area $22$, we can use \href{https://artofproblemsolving.com/wiki/index.php?title=Pick\%27s\_Theorem}{Pick's Theorem} to solve for $I$ as follows:\begin{align*}I + \frac B 2 -1 &= 22 \\I + \frac{24}{2} &= 23 \\I + 12 &= 23 \\I &= 11\end{align*}
Because the total amount of lattice points inside the polygon (and, consequently, inside the shaded region) is $I+B$, our answer is $11+24=\boxed{\textbf{(B) }35}$.

 


	\item (\textit{AHSME 1957}) In $\triangle ABC, AC = CD$ and $\angle CAB - \angle ABC = 30^\circ$. Then $\angle BAD$ is: 


 
                    \begin{figure}[H]
                        \centering
                        \includegraphics[width=0.5\textwidth]{Problem/AHSME_I_1957_44_Problem_Figure_1.png}
                    \end{figure}
                    
$\textbf{(A)}\ 30^\circ\qquad\textbf{(B)}\ 20^\circ\qquad\textbf{(C)}\ 22\frac{1}{2}^\circ\qquad\textbf{(D)}\ 10^\circ\qquad\textbf{(E)}\ 15^\circ$

 \textbf{Answer: }E



\textbf{Solution}

Because $\triangle ACD$ is \href{https://artofproblemsolving.com/wiki/index.php?title=Isosceles\_triangle}{isosceles} with $AC=CD$, $\measuredangle CAD=\measuredangle CDA=\theta$, where $\theta$ is some angle measure. Because $\angle CDA$ and $\angle ADB$ form a \href{https://artofproblemsolving.com/wiki/index.php?title=Straight\_angle}{straight angle}, $\measuredangle ADB = 180^{\circ}-\theta$. Thus, because the interior angles of a triangle add to $180^{\circ}$, $\measuredangle BAD + \measuredangle ABC = \theta$, so $\measuredangle ABC = \theta - \measuredangle BAD$. Notice that $\measuredangle CAB = \theta + \measuredangle BAD$. With all of this information and recalling that, from the problem, $\measuredangle CAB - \measuredangle ABC = 30^{\circ}$, we see that:\begin{align*}\measuredangle CAB - \measuredangle ABC &= 30^{\circ} \\\theta + \measuredangle BAD - (\theta - \measuredangle BAD) &= 30^{\circ} \\2\measuredangle BAD &= 30^{\circ} \\\measuredangle BAD &= 15^{\circ}\end{align*}
Thus, our answer is $\boxed{\textbf{(E) }15^{\circ}}$.

 


	\item (\textit{AHSME 1957}) If two real numbers $x$ and $y$ satisfy the equation $\frac{x}{y} = x - y$, then: 


 $\textbf{(A)}\ {x \ge 4}\text{ or }{x \le 0}\qquad \\  \textbf{(B)}\ {y}\text{ can equal }{1}\qquad \\  \textbf{(C)}\ \text{both }{x}\text{ and }{y}\text{ must be irrational}\qquad\\  \textbf{(D)}\ {x}\text{ and }{y}\text{ cannot both be integers}\qquad\\  \textbf{(E)}\ \text{both }{x}\text{ and }{y}\text{ must be rational}$

 \textbf{Answer: }A



\textbf{Solution}

From the initial equation $\frac x y = x-y$, we can solve for $y$ in terms of $x$ as follows:\begin{align*}\\frac x y &= x-y \\x &= xy-y^2 \\y^2-xy+x &= 0 \\y = \frac{x \pm \sqrt{x^2-4x}}2\end{align*}
Because $y$ is real, the \href{https://artofproblemsolving.com/wiki/index.php?title=Discriminant}{discriminant} of the above expression for $y$ must be $\geq 0$, so $x^2-4x=x(x-4) \geq 0$. This is true when $x \geq 4$ and $x \leq 0$, so we choose answer $\boxed{\textbf{(A)}}$.

 


	\item (\textit{AHSME 1957}) Two perpendicular chords intersect in a circle. The segments of one chord are $3$ and $4$; 
the segments of the other are $6$ and $2$. Then the diameter of the circle is: 


 $\textbf{(A)}\ \sqrt{89}\qquad  \textbf{(B)}\ \sqrt{56}\qquad  \textbf{(C)}\ \sqrt{61}\qquad  \textbf{(D)}\ \sqrt{75}\qquad \textbf{(E)}\ \sqrt{65}$

 \textbf{Answer: }E



\textbf{Solution 1}


                    \begin{figure}[H]
                        \centering
                        \includegraphics[width=0.5\textwidth]{Solution/AHSME_I_1957_46_Solution_Figure_1.png}
                    \end{figure}
                    

 Let the chords intersect the circle at points $A,B,C,$ and $D$ to form simple polygon $ABCD$. Further, let the chords intersect at point $P$ with $AP=6,PC=2,BP=3,$ and $PD=4$, as in the diagram. Then, because $\overline{AC} \perp \overline{BD}$, by the \href{https://artofproblemsolving.com/wiki/index.php?title=Pythagorean\_Theorem}{Pythagorean Theorem}, $AB=3\sqrt5$ and $BC=\sqrt{13}$. Because a circle is determined by three coplanar points, the \href{https://artofproblemsolving.com/wiki/index.php?title=Circumcircle}{circumcircle} of $\triangle ABC$ will be the circumcircle of $ABCD$, so the circumdiameter of $\triangle ABC$ will be our desired answer. We know that $[\triangle ABC]=\frac{(6+2) \cdot 3}2=12$. Furthermore, we know that we can express this area as $\frac{abc}{4R}$, where $a,b,$ and $c$ are $\triangle ABC$'s side lengths and $R$ is its circumradius. Setting this expression equal to $12$, we can now solve for $R$:\begin{align*}\frac{abc}{4R} &= 12 \\\frac{\sqrt{13} \cdot 8 \cdot 3\sqrt{15}}{4R} &= 12 \\\frac{2 \cdot 3\sqrt{65}}R &= 12 \\6\sqrt{65} &= 12R \\R &= \frac{\sqrt{65}} 2\end{align*}
Because the problem asks for the diameter of the circle, our answer is $2R=2 \cdot \frac{\sqrt{65}}2=\boxed{\textbf{(E) }\sqrt{65}}$.

 



\textbf{Solution 2 (Answer choices, intution)}

Because one of the chords decribed in the problem has length $6+2=8$, we know that the diameter, the largest chord in the circle, must have a length $\geq 8$. This fact eliminates options (B) and (C). Furthermore, because the chord of length $8$ is close to bisecting the other perpendicular chord (being only $0.5$ off of its midpoint), it should be rather close to the length of the diameter. Because the circle's tangent lines are perpendicular to the diameter, moving a small distance away from the diameter will not rapidly decrease the length of the chord, so $AC$ is only slightly less than the diameter's length. This fact guides us towards answer $\boxed{\textbf{(E) }\sqrt{65}}$.

 


	\item (\textit{AHSME 1957}) In circle $O$, the midpoint of radius $OX$ is $Q$; at $Q$, $\overline{AB} \perp \overline{XY}$. The semi-circle with $\overline{AB}$ 
as diameter intersects $\overline{XY}$ in $M$. Line $\overline{AM}$ intersects circle $O$ in $C$, and line $\overline{BM}$ 
intersects circle $O$ in $D$. Line $\overline{AD}$ is drawn. Then, if the radius of circle $O$ is $r$, $AD$ is: 


 
                    \begin{figure}[H]
                        \centering
                        \includegraphics[width=0.5\textwidth]{Problem/AHSME_I_1957_47_Problem_Figure_1.png}
                    \end{figure}
                    
$\textbf{(A)}\ r\sqrt{2}\qquad\textbf{(B)}\ r\qquad\textbf{(C)}\ \text{not a side of an inscribed regular polygon}\qquad\textbf{(D)}\ \frac{r\sqrt{3}}{2}\qquad\textbf{(E)}\ r\sqrt{3}$

 \textbf{Answer: }A



\textbf{Solution}

Because $\overline{XY}$ is a diameter of the circle and $\overline{AB} \perp \overline{XY}$, we know that $\overline{XY}$ bisects $\overline{AB}$, so $AQ=BQ$. Thus, $M$ is on the \href{https://artofproblemsolving.com/wiki/index.php?title=Perpendicular\_bisector\#Locus}{perpendicular bisector} of $\overline{AB}$, and so $AM=BM$. Furthermore, by \href{https://artofproblemsolving.com/wiki/index.php?title=Thales\%27\_Theorem}{Thales' Theorem}, $\measuredangle AMB = 90^{\circ}$. Thus, because $\triangle AMB$ is a right isosceles triangle with $AM=BM$, it is a \href{https://artofproblemsolving.com/wiki/index.php?title=45-45-90\_triangle}{45-45-90 triangle}. Thus, $\measuredangle ABM = 45^{\circ}$. Now, draw $\overline{OA}$ and $\overline{OD}$. Because $\angle ABD \cong \angle ABM$ is an \href{https://artofproblemsolving.com/wiki/index.php?title=Inscribed\_angle}{inscribed angle} which intercepts minor arc $AD$, the measure of \href{https://artofproblemsolving.com/wiki/index.php?title=Central\_angle}{central angle} $\angle AOD$ must be $2\measuredangle ABD = 2 \cdot 45^{\circ} = 90^{\circ}$. Because $OA=OD=r$, $\triangle AOD$ is also a $45$-$45$-$90$ triangle, so $AD=\boxed{\textbf{(A) }r\sqrt2}$.

 


	\item (\textit{AHSME 1957}) Let $ABC$ be an equilateral triangle inscribed in circle $O$. $M$ is a point on arc $BC$. 
Lines $\overline{AM}$, $\overline{BM}$, and $\overline{CM}$ are drawn. Then $AM$ is: 


 
                    \begin{figure}[H]
                        \centering
                        \includegraphics[width=0.5\textwidth]{Problem/AHSME_I_1957_48_Problem_Figure_1.png}
                    \end{figure}
                    
$\textbf{(A)}\ \text{equal to }{BM + CM}\qquad  \textbf{(B)}\ \text{less than }{BM + CM}\qquad \\ \textbf{(C)}\ \text{greater than }{BM+CM}\qquad \\ \textbf{(D)}\ \text{equal, less than, or greater than }{BM + CM}\text{, depending upon the position of } {M}\qquad  \\ \textbf{(E)}\ \text{none of these}$

 \textbf{Answer: }A



\textbf{Solution}

Since quadrilateral $ABMC$ is inscribed in circle $O$, thus it is a cyclic quadrilateral. By \href{https://artofproblemsolving.com/wiki/index.php?title=Ptolemy\%27s\_Theorem}{Ptolemy's Theorem}, 
 \[AC \cdot MB + MC \cdot AB = BC \cdot AM.\] Because $\triangle ABC$ is equilateral, we cancel out $AB$, $AC$, and $BC$ to get that 
 \[BM + CM = AM \implies \boxed{\textbf{(A)}}.\]



	\item (\textit{AHSME 1957}) The parallel sides of a trapezoid are $3$ and $9$. The non-parallel sides are $4$ and $6$. 
A line parallel to the bases divides the trapezoid into two trapezoids of equal perimeters. 
The ratio in which each of the non-parallel sides is divided is: 


 
                    \begin{figure}[H]
                        \centering
                        \includegraphics[width=0.5\textwidth]{Problem/AHSME_I_1957_49_Problem_Figure_1.png}
                    \end{figure}
                    
$\textbf{(A)}\ 4: 3\qquad\textbf{(B)}\ 3: 2\qquad\textbf{(C)}\ 4: 1\qquad\textbf{(D)}\ 3: 1\qquad\textbf{(E)}\ 6: 1$

 \textbf{Answer: }C



\textbf{Solution}


                    \begin{figure}[H]
                        \centering
                        \includegraphics[width=0.5\textwidth]{Solution/AHSME_I_1957_49_Solution_Figure_1.png}
                    \end{figure}
                    

 Let the points be labeled as in the new diagram above, with $BF=6$ and $CG=4$ (from the problem). Because $\overline{BC} \parallel \overline{FG}$ and $\frac{BC}{FG}=\frac 1 3$, $AB=\frac{AF}3=\frac{AB+6}3$ and $AC=\frac{AG}3=\frac{AC+4}3$. Solving these equations for $AB$ and $AC$, respectively yields $AB=3$ and $AC=2$. Let $DE=x$. Thus, because $\overline{DE} \parallel \overline{BC}$, $\frac{BC}{DE}=\frac 3 x=\frac{AB}{AD}=\frac3{3+BD}$. Solving this equation for $BD$ yields $BD=x-3$. Similarly, $\frac 3 x=\frac{AC}{AE}=\frac2{2+CE}$, and solving this equation for $CE$ yields $CE=\frac{2x}3-2$.

 Now, we can set the perimeters of $BCED$ and $DEGF$ equal to each other to solve for $x$:\begin{align*}3+(x-3)+x+(\frac{2x}3-2) &= x+[6-(x-3)]+9+[4-(\frac{2x}3-2)] \\\frac{8x}3-2 &= 9+9+4-\frac{2x}3+2 \\\frac{10x}3 &= 26 \\\frac{5x}3 &= 13 \\5x &= 39 \\x &= \frac{39}5\end{align*}
To find the ratio $\frac{BD}{DF}=\frac{x-3}{6-(x-3)}$, we substitute $x=\frac{39}5$ into this expression to find our answer:\begin{align*}\frac{x-3}{6-(x-3)} &= \frac{x-3}{9-x} \\&= \frac{\frac{39-15}5}{\frac{45-39}5} \\&= \frac{24}6 \\&= \frac 4 1\end{align*}
Thus, our answer is $\boxed{\textbf{(C) }4:1}$.

 


	\item (\textit{AHSME 1957}) In circle $O$, $G$ is a moving point on diameter $\overline{AB}$. $\overline{AA'}$ is drawn perpendicular to $\overline{AB}$ and equal to $\overline{AG}$. 
$\overline{BB'}$ is drawn perpendicular to $\overline{AB}$, on the same side of diameter $\overline{AB}$ as $\overline{AA'}$, 
and equal to $BG$. Let $O'$ be the midpoint of $\overline{A'B'}$. Then, as $G$ moves from $A$ to $B$, point $O'$: 


 $\textbf{(A)}\ \text{moves on a straight line parallel to }{AB}\qquad \\ \textbf{(B)}\ \text{remains stationary}\qquad\\  \textbf{(C)}\ \text{moves on a straight line perpendicular to }{AB}\qquad\\  \textbf{(D)}\ \text{moves in a small circle intersecting the given circle}\qquad\\  \textbf{(E)}\ \text{follows a path which is neither a circle nor a straight line}$

 \textbf{Answer: }B



\textbf{Solution}


                    \begin{figure}[H]
                        \centering
                        \includegraphics[width=0.5\textwidth]{Solution/AHSME_I_1957_50_Solution_Figure_1.png}
                    \end{figure}
                    

 Let $AG=AA'=a$ and $BG=BB'=b$. Then, we know that the diameter $d$ of the circle equals $AG+BG=a+b$. Thus, because the circle's diameter does not change, $a+b$ is constant. 

 Because $A'O'=O'B'$ and $AO=OB$, $\overline{OO'} \parallel \overline{AA'}$. Thus, $\overline{OO'} \perp \overline{AB}$, and so $OO'$ is the distance from $O'$ to $\overline{AB}$. 

 Let $P$ be some point which moves along $\overline{A'B'}$. Because $A'B'$ is a line segment, as $P$ moves from $A'$ to $B'$, its distance from $\overline{AB}$ will vary linearly with how much it has travelled along $\overline{A'B'}$. Thus, when it is halfway along $\overline{A'B'}$ (in other words, when $P=O'$) its distance from $\overline{AB}$ will be the arithmetic mean of its distance from $\overline{AB}$ at $A'$ (namely, $A'A=a$) and its distance from $\overline{AB}$ at $B'$ (namely, $B'B=b$). Thus, $O'O = \frac{a+b}2$. 

 Because $a+b=d$, a constant, $\frac{a+b}2$ is a constant as well. Thus, $OO'$ is the same regardless of the position of $G$. Furthermore, from our work in paragraph 2, we know that $O'$ must lie on the line perpendicular to $\overline{AB}$ through point $O$. Therefore, because $O'$ is a fixed distance from a fixed point on a fixed line, and it will not suddenly "jump across" to the other side of $\overline{AB}$, we can say with confidence that point $O'$ $\fbox{\textbf{(B)} remains stationary}$.

 


\end{enumerate}